% =========================================================================
% PROJETO DE PESQUISA (PIBIC / CNPQ)
% =========================================================================
\documentclass{amsart}
\usepackage[a4paper, total={6in, 8in}]{geometry}

\usepackage{amssymb}
\usepackage{amsmath}
\usepackage{amsthm}
\usepackage{bm}
\usepackage{enumitem}
\usepackage{mathtools}
\usepackage[brazilian]{babel}
\usepackage[utf8]{inputenc}
\usepackage[T1]{fontenc}
\usepackage[pdftex]{graphicx}
\usepackage{color}
\usepackage{tikz}
\usetikzlibrary{arrows, calc}
\newtheorem{prop}{Proposição}
\usepackage{verbatim}
\usepackage[hyphens]{url}
\usepackage{longtable}

\newcommand{\numedital}{@@EDITAL@@}
\newcommand{\longtitle}{@@RESEARCH-TITLE@@: @@RESEARCH-SUBTITLE@@}

\usepackage{setspace}
\doublespacing

\begin{document}
\title[@@RESEARCH-TITLE@@]{\longtitle}
\maketitle
\thispagestyle{empty}
\vfill

\begin{flushright}
\begin{minipage}{8cm}
\begin{flushright}
Projeto de Iniciação Científica\\submetido para avaliação no Edital: \numedital
\end{flushright}
\end{minipage}
\end{flushright}

\vfill

\renewcommand{\arraystretch}{2}
\begin{tabular}{ll}
\textbf{Título do projeto:} &
\begin{minipage}[t]{8cm}
\longtitle
\end{minipage}
\\
\textbf{Palavras-chave do projeto:} &
\begin{minipage}[t]{8cm}
Otimização Combinatória, Algoritmos, C++23
\end{minipage}\\
\textbf{Área do conhecimento do projeto:} &
\begin{minipage}[t]{8cm}
Ciência da Computação, Matemática Discreta
\end{minipage}
\end{tabular}
\newpage

\tableofcontents
\newpage

% -------------------------------------------------------------------------
\section{Resumo}
% -------------------------------------------------------------------------
@@RESEARCH-DESC@@

% -------------------------------------------------------------------------
\section{Introdução e Justificativa}
% -------------------------------------------------------------------------
Apresente a contextualização científica do problema, as motivações matemáticas e computacionais e as principais referências da área~\cite{cormen2022introduction, korte2018combinatorial}.

% -------------------------------------------------------------------------
\section{Objetivos}
% -------------------------------------------------------------------------
\subsection{Objetivo geral}
Proporcionar formação sólida em @@RESEARCH-TITLE@@, integrando teoria formal, implementação eficiente em C++23 e validação experimental.

\subsection{Objetivos específicos}
\begin{itemize}
    \item Estudo dos fundamentos matemáticos e garantias estruturais do domínio;
    \item Implementação orientada a objetos e modular em C++23;
    \item Execução de experimentos computacionais controlados e análise comparativa.
\end{itemize}

% -------------------------------------------------------------------------
\section{Metodologia}
% -------------------------------------------------------------------------
A metodologia combina ciclos iterativos de estudo teórico, desenvolvimento em C++23 sob interface polimórfica comum e avaliação experimental reprodutível.

% -------------------------------------------------------------------------
\section{Cronograma de Atividades}
% -------------------------------------------------------------------------
\begin{longtable}{|c|c|c|c|c|c|c|c|c|c|c|c|c|}
\hline
Atividade   &01&02&03&04&05&06&07&08&09&10&11&12\\ \hline
\endfirsthead
\hline
Atividade   &01&02&03&04&05&06&07&08&09&10&11&12\\ \hline
\endhead
Ativ. 1--3 &\textbf{\textsf{X}}&\textbf{\textsf{X}}&\textbf{\textsf{X}}&\textbf{\textsf{X}}&&&&&&&&\\ \hline
Ativ. 4--6 &&&&\textbf{\textsf{X}}&\textbf{\textsf{X}}&\textbf{\textsf{X}}&&&&&&\\ \hline
Ativ. 7--9 &&&&&&&\textbf{\textsf{X}}&\textbf{\textsf{X}}&\textbf{\textsf{X}}&\textbf{\textsf{X}}&&\\ \hline
Ativ. 10--11 &&&&&&&&&&\textbf{\textsf{X}}&\textbf{\textsf{X}}&\textbf{\textsf{X}}\\ \hline
\end{longtable}

% -------------------------------------------------------------------------
\section{Viabilidade}
% -------------------------------------------------------------------------
O projeto requer apenas estação de trabalho convencional e ferramentas abertas (compilador C++23, Python e distribuição LaTeX).

\newpage
\bibliographystyle{plain}
\bibliography{common/references}

\end{document}
